
\section{Conclusion}
\label{sec:conclusion}
We have presented \KFOL, a symmetric $K$-player calculus with a minimal,
prenex core and a $K=2$ specialization that is algebraically rich and
serves as a retract of classical \FOL. By internalizing player roles,
vectorized outcomes, and explicit $S_K$ symmetry, the calculus
provides a foundation for analyzing multi-agent games. We extended this
core with IF-style binders, which raises the complexity of the
forceability problem from PSPACE-complete to NEXPTIME-hard,
and allows the logic to syntactically capture information constraints.

We want to be explicit about the limits of what this paper shows. In
particular, the complexity results should \emph{not} be read as evidence that
$K>2$ roles carry strategic content beyond the two-player machinery the
calculus is built from: as discussed in \Cref{sec:k-complexity}, the
PSPACE-completeness of the broadcast core is uniform in $K$, coalitional
forceability under broadcast visibility reduces to ordinary two-role
forceability of a merged game (\Cref{prop:coalition-merge-lossless}), and even
the NEXPTIME-hardness result is witnessed by a $K=2$ construction
(\Cref{rem:nexptime-is-k2}). What we have established is that a specific
$K$-role, prenex, first-order game syntax admits clean semantics, a retract
onto classical \FOL at $K=2$, a nontrivial (if only partially understood)
symmetry algebra, and MPST-style protocol readings, with the expected
metatheory worked out. We have \emph{not} established that this syntax does
anything a family of two-player games could not already do; that is the open
question we consider most important, first below.

\paragraph{Open questions.}
\begin{itemize}
  \item \textbf{Genuine $K\ge 3$ separation.} Is there a \KFOL\ phenomenon ---
  in forceability, complexity, or elsewhere --- that is not reducible to a
  family of two-player or two-coalition facts? \Cref{subsec:coalitional-forceability}
  conjectures that coalitional forceability under IF-binders, where a
  coalition's own members are mutually information-restricted \emph{and}
  whose objectives are coupled across that restriction, is the place to look:
  role-merging is provably lossless once visibility is broadcast
  (\Cref{prop:coalition-merge-lossless}), a mere $\setminus V$-crossing within
  a coalition is not by itself enough
  (\Cref{prop:cross-diagonal-forceability}), and the existing NEXPTIME-hardness
  construction does not touch this case (\Cref{rem:nexptime-is-k2}). No
  separating construction is currently known; finding one --- or proving that
  none exists, i.e.\ that \KFOL\ forceability always reduces to two-coalition
  facts --- would be the paper's missing central theorem. This same fork has
  now recurred three times under three different readings of the calculus ---
  here, in the tensor/coalition-merging discussion of
  \Cref{subsec:open-games}, and in the Byzantine-robustness remark of
  \Cref{sec:framework} --- and collapsed to two-coalition facts in the
  perfect-information case each time. That convergence is itself evidence
  about where to look: every attempt so far suggests genuine $K\ge 3$ content,
  if it exists in this calculus, comes from information asymmetry between
  roles rather than from coalition counting. All of this is posed
  information-theoretically; \Cref{subsec:computational-soundness-gap}
  poses, but does not resolve, the same question once independence is only
  assumed to hold computationally.

  \item \textbf{A third complexity regime between the two known ones ---
  and a more basic gap underneath it.} \Cref{rem:hierarchical-if} separates
  three claims of different strength. The most basic is not about the
  hierarchical restriction at all: \Cref{prop:nexptime-if} only proves
  NEXPTIME-\emph{hardness} of the general IF fragment, with no matching
  upper bound, so unrestricted \KFOL\ forceability is not currently known to
  be decidable --- Peterson--Reif--Azhar's own finding that fully general
  (non-hierarchical) multiplayer incomplete-information games can be
  undecidable~\cite{PetersonReifAzhar2001LowerBounds,PetersonReifAzhar2002DecisionAlgorithms}
  leaves this as a live possibility, not merely a formality. Restricting to
  \emph{hierarchical} (pairwise-comparable) $\setminus V$ patterns is very
  likely decidable by a routine Skolem-table argument (not a deep unknown);
  what is genuinely open is whether its complexity matches their
  $(d{+}1)$-fold-exponential bound exactly, with a matching lower bound
  specific to \KFOL's own syntax. This is a well-posed, independently
  interesting complexity question even
  if the $K\ge 3$ separation conjecture above turns out to be false, and a
  natural warm-up for it: it isolates nesting depth as a complexity
  parameter distinct from $K$, and --- because
  \Cref{prop:cross-diagonal-forceability} already shows comparability alone
  does not break flattening --- confirms that a separating construction, if
  one exists, must specifically use non-hierarchical independence, not
  merely nontrivial independence.

  \item \textbf{Relationship to ATL, Coalition Logic, and Strategy Logic ---
  check this one first.} Individual forceability $F_i$ is semantically a
  one-shot instance of ATL's $\langle\!\langle i \rangle\!\rangle$ modality
  (\Cref{sec:related}). We have not worked out a translation in either
  direction, nor compared complexity results on equal footing (formula
  complexity here vs.\ combined complexity there). Unlike the other items in
  this list, a bad answer here is not merely a missing result but a threat to
  the rest of them: if forceability turns out to be a routine encoding of
  $\langle\!\langle i \rangle\!\rangle$, several other open questions in this
  list inherit that reduction for free (in either a good or a bad way,
  depending on which direction it runs) and should be revisited in its light.
  It is also comparatively cheap to check. This should be done before any
  claim of novelty relative to that literature is treated as established
  rather than plausible, and before investing in the harder items below.

  \item \textbf{Team semantics / Dependence Logic.} The IF-binder extension
  overlaps with Väänänen's dependence logic~\cite{Vaananen2007DependenceLogic},
  the modern reformulation of the same two-player, imperfect-information
  territory Hintikka--Sandu IF logic occupies. We conjectured team semantics
  corresponds to the $K=2$ slice of \KFOL's information-sensitive fragment,
  consistent with \Cref{rem:nexptime-is-k2}; on reflection a clean
  biconditional is unlikely to be the right target, since IF logic and
  dependence logic are already known to diverge at the open-formula level
  (despite sentence-level equivalence) and team semantics itself has
  lax/strict variants. The more realistic --- and still worthwhile --- shape
  of this item is a \emph{taxonomy}: which variant corresponds to \KFOL's
  $K=2$ IF-fragment under which translation, and exactly where the others
  diverge from it. We have not attempted this in either form.

  \item \textbf{A full proof theory --- scope this before attempting it.} We
  have no sequent calculus, cut rule, or cut-elimination result for \KFOL;
  \Cref{sec:ra-view}'s relational-algebra reading gives compositional
  intuition (profiles as joins) but is not a substitute for one. This should
  not be stated as flatly open, though: classical IF logic is provably
  \emph{not} axiomatizable (validity is not recursively enumerable, via
  Trakhtenbrot's theorem~\cite{Trakhtenbrot1950,sep-logic-if}), and our
  IF-binder extension is close enough to IF logic
  (\Cref{subsec:if-binders,sec:related}) that a complete proof calculus for
  the \emph{full IF-extended} fragment may be impossible in principle, not
  merely unbuilt. What remains open --- and what a proof theory should
  actually be scoped to --- is the broadcast core (which retracts onto
  classical \FOL at $K=2$ and inherits its axiomatizability) and whichever
  restricted IF fragments (if any) stay decidable or recursively
  axiomatizable; the conjecture that the full IF-extended fragment resists
  complete axiomatization should itself be checked, not assumed, since we
  have not verified \KFOL's IF-binders are expressive enough to actually
  embed the constructions Trakhtenbrot's argument needs.

  \item \textbf{\KFOL\ as a semantic model of a $K$-player fragment of linear
  logic.} This would mean an abstraction from multi-party linear propositions
  into \KFOL\ games such that cut corresponds to game composition. Carbone,
  Montesi, Sch\"urmann, and Yoshida's coherence-proofs correspondence for
  MPST~\cite{CarboneMontesiSchurmannYoshida2015Coherence} is the closest
  existing result along these lines we are aware of, though it is fully
  cooperative (coherence, not our adversarial forceability); an adversarial
  analogue of coherence is, as far as we know, still to be worked out
  (\Cref{sec:related}).

  \item \textbf{Classifying the complement operator.}
  \Cref{rem:no-classification} notes that our negative results about
  $\CompOp$ (failure to preserve zero-sumness, failure of De Morgan duality)
  are witnessed by single counterexamples, not a classification of when these
  properties hold. Relatedly (\Cref{rem:outcome-geometry}): what group of
  transformations of the $K$-player outcome space interpolates between the
  one-hot (zero-sum) and diagonal (fully cooperative) extremes for $K>2$,
  generalizing the $K=2$ rotation/complement symmetries of
  \Cref{sec:symmetry}, and do its orbits correspond to meaningful degrees of
  cooperativeness or coalition structure? A categorical axiomatization of what
  the group action preserves (in the spirit of $*$-autonomous categories for
  classical duality) would likely answer both questions at once. A more
  speculative version of the same question, not derived here: the $S_K$
  action also acts on coalition families
  $\mathcal S\subseteq 2^{\{0,\dots,K-1\}}$
  (\Cref{rem:coalition-operator-parameters}), and its orbits are a candidate
  notion of ``type of power structure'' --- ``no singleton dominates'' is
  $S_K$-invariant, ``role 0 is hegemon'' is not --- which would make
  ``unbiased'' formal as $S_K$-equivariance of a coalition predicate. We have
  not computed a single orbit or posed the corresponding representation
  question (which $S_K$-invariant labelings of the coalition lattice are
  realized by an actual \KFOL\ formula's forceability landscape); this is a
  research direction, not a claim.

  \item \textbf{Compositional game theory as candidate machinery.}
  \Cref{subsec:open-games} surveys open games, lenses/optics, and
  Escard\'o--Oliva selection functions as a possible source for both the
  categorical axiomatization above and a precise version of the
  coalitional-merging conjecture (do coalitions correspond to lenses, with
  the over-/under-approximation gap identifying exactly where $K\ge3$ content
  lives?). One specific obstacle: the standard open-games equilibrium functor
  is Nash-style, which is the composability failure mode \KFOL's forceability
  semantics was built to avoid, so this machinery cannot be adopted wholesale
  without first substituting an adversarial/minimax selection notion ---
  itself apparently unstudied in that literature. This is scoped as a side
  investigation to try, not a claim about what it will find.

  \item \textbf{Probabilistic and quantitative extensions.} Combining this
  work with mixed strategies, graded payoffs, or probabilistic reasoning is
  future work; \Cref{subsec:payoff-ignorance} notes specifically that the
  $\{0,1\}$-valued setting collapses the utility-gap arguments used in
  rational-proof and scoring-rule literature, which graded payoffs would
  restore.
\end{itemize}
